% ============================================================================
%  ANÁLISIS DE PRUEBAS — ChibchaWeb
%  Documento compilable con pdflatex o lualatex
% ============================================================================
\documentclass[12pt,a4paper]{article}

% — Codificación y tipografía ------------------------------------------------
\usepackage[utf8]{inputenc}
\usepackage[T1]{fontenc}
\usepackage[spanish]{babel}
\usepackage{lmodern}

% — Diseño de página ---------------------------------------------------------
\usepackage[margin=2.5cm]{geometry}
\usepackage{parskip}

% — Tablas y figuras ----------------------------------------------------------
\usepackage{booktabs}
\usepackage{tabularx}
\usepackage{longtable}
\usepackage{multirow}
\usepackage{array}

% — Colores y resaltado -------------------------------------------------------
\usepackage[table,dvipsnames]{xcolor}
\definecolor{aprobado}{HTML}{27AE60}
\definecolor{aprobadotodo}{HTML}{F39C12}
\definecolor{reprobado}{HTML}{E74C3C}
\definecolor{sinprueba}{HTML}{95A5A6}
\definecolor{codebg}{HTML}{F5F5F5}
\definecolor{secbg}{HTML}{EBF5FB}

% — Listados de código --------------------------------------------------------
\usepackage{listings}
\lstset{
  backgroundcolor=\color{codebg},
  basicstyle=\ttfamily\footnotesize,
  breaklines=true,
  frame=single,
  rulecolor=\color{gray!40},
  numbers=left,
  numberstyle=\tiny\color{gray},
  numbersep=6pt,
  xleftmargin=1.5em,
  framexleftmargin=1.5em,
  captionpos=b,
  tabsize=4,
  showstringspaces=false,
  literate={á}{{\'a}}1 {é}{{\'e}}1 {í}{{\'i}}1 {ó}{{\'o}}1 {ú}{{\'u}}1
           {ñ}{{\~n}}1 {Á}{{\'A}}1 {É}{{\'E}}1 {Í}{{\'I}}1 {Ó}{{\'O}}1
           {Ú}{{\'U}}1 {Ñ}{{\~N}}1,
}

% — Hipervínculos y metadatos -------------------------------------------------
\usepackage{hyperref}
\hypersetup{
  colorlinks=true,
  linkcolor=MidnightBlue,
  urlcolor=MidnightBlue,
  pdfauthor={Equipo ChibchaWeb},
  pdftitle={Análisis de pruebas — ChibchaWeb},
}

% — Encabezados y pies de página ----------------------------------------------
\usepackage{fancyhdr}
\pagestyle{fancy}
\fancyhf{}
\fancyhead[L]{\small\textsc{ChibchaWeb}}
\fancyhead[R]{\small Análisis de pruebas}
\fancyfoot[C]{\thepage}
\renewcommand{\headrulewidth}{0.4pt}

% — Comandos auxiliares --------------------------------------------------------
\newcommand{\aprobado}{\textcolor{aprobado}{\textbf{APROBADO}}}
\newcommand{\aprobadotodo}{\textcolor{aprobadotodo}{\textbf{APROBADO (TODO)}}}
\newcommand{\aprobadojust}{\textcolor{aprobadotodo}{\textbf{APROBADO*}}}
\newcommand{\reprobado}{\textcolor{reprobado}{\textbf{REPROBADO}}}
\newcommand{\sinprueba}{\textcolor{sinprueba}{\textbf{SIN PRUEBAS}}}

% ============================================================================
\begin{document}

% — Portada -------------------------------------------------------------------
\begin{titlepage}
\centering
\vspace*{3cm}
{\Huge\bfseries Análisis de Pruebas\\[0.4em]
Proyecto ChibchaWeb\par}
\vspace{1.5cm}
{\Large\itshape Informe de ejecución y clasificación\\de la suite de pruebas de rendimiento\par}
\vspace{2cm}
{\large
\begin{tabular}{rl}
\textbf{Proyecto:}   & ChibchaWeb (Django 5.2.4) \\
\textbf{Base de datos:} & SQLite 3 (desarrollo) \\
\textbf{Framework de pruebas:} & pytest 9.0.3 + pytest-django 4.12.0 \\
\textbf{Entorno:}    & Python 3.10.11 — Windows 10 \\
\textbf{Fecha:}      & 27 de mayo de 2026 \\
\end{tabular}\par}
\vfill
{\small Documento generado automáticamente a partir de la ejecución de la suite de pruebas.}
\end{titlepage}

% — Índice --------------------------------------------------------------------
\tableofcontents
\newpage

% ============================================================================
\section{Introducción}
% ============================================================================

El presente documento contiene el análisis detallado de las pruebas ejecutadas sobre el
proyecto \textbf{ChibchaWeb}, una plataforma web desarrollada con Django~5.2.4 para la
gestión de dominios, clientes, pagos y soporte técnico.

Se ejecutaron \textbf{10 pruebas} utilizando \texttt{pytest} con la base de datos de
desarrollo (SQLite). Todas las pruebas fueron aprobadas exitosamente.

\subsection{Resumen ejecutivo}

\begin{center}
\rowcolors{2}{white}{secbg}
\begin{tabular}{lc}
\toprule
\textbf{Métrica}          & \textbf{Valor} \\
\midrule
Total de pruebas          & 10 \\
Pruebas aprobadas         & 10 \\
Pruebas reprobadas        & 0  \\
Errores de ejecución      & 0  \\
Advertencias (warnings)   & 10 (marcas no registradas) \\
Tiempo total de ejecución & 38.33 s \\
Código de salida          & 0 (éxito) \\
\bottomrule
\end{tabular}
\end{center}

\subsection{Clasificación de pruebas}

Las pruebas se clasifican en cinco categorías según el tipo de verificación que realizan:

\begin{enumerate}
  \item \textbf{Unitarias:} verifican componentes individuales aislados (capa ORM).
  \item \textbf{Integración:} verifican la interacción entre componentes (vistas HTTP + ORM + servicios simulados).
  \item \textbf{Sistema:} verifican flujos completos de extremo a extremo.
  \item \textbf{Eficiencia:} verifican tiempos de respuesta, rendimiento y uso de recursos.
  \item \textbf{Robustez:} verifican estabilidad bajo carga, concurrencia y condiciones adversas.
\end{enumerate}

% ============================================================================
\newpage
\section{Pruebas Unitarias}
% ============================================================================

Las pruebas unitarias verifican el correcto funcionamiento de componentes individuales
de forma aislada. En esta suite, se identifican dos pruebas que operan exclusivamente
sobre la capa de datos (ORM de Django) sin involucrar vistas HTTP ni servicios externos.

\subsection{Resultados}

\begin{center}
\rowcolors{2}{white}{secbg}
\begin{tabularx}{\textwidth}{>{\raggedright}p{3.8cm} >{\raggedright}p{3cm} r r >{\centering\arraybackslash}p{2.5cm}}
\toprule
\textbf{Prueba} & \textbf{Módulo} & \textbf{Tiempo (s)} & \textbf{Umbral (s)} & \textbf{Resultado} \\
\midrule
Búsqueda/filtro de dominios (ORM)
  & Dominios ORM
  & 0.002868
  & $< 1.0$
  & \aprobadotodo \\
Actualización de dominio (ORM)
  & Dominios ORM
  & 0.000862
  & $< 3.0$
  & \aprobadotodo \\
\bottomrule
\end{tabularx}
\end{center}

\subsection{Descripción de las pruebas}

\subsubsection{Búsqueda y filtro de dominios}

\begin{itemize}
  \item \textbf{Archivo:} \texttt{test\_domain\_search\_filter\_performance.py}
  \item \textbf{Escenario:} Se crean 200 dominios con prefijo \texttt{qa-search} y 50 dominios con
        prefijo \texttt{otro} (compra distribuidor). Se ejecuta una consulta ORM con filtros
        \texttt{icontains} y \texttt{compraDistribuidor=False}, limitada a 25 resultados ordenados.
  \item \textbf{Verificación:} Se comprueba que se obtienen exactamente 25 resultados y que el
        tiempo de ejecución es inferior a 1.0~s.
  \item \textbf{Nota:} Marcada con TODO porque falta una vista/URL de búsqueda en el módulo
        \texttt{Dominios} para validar el flujo HTTP completo.
\end{itemize}

\subsubsection{Actualización de dominio}

\begin{itemize}
  \item \textbf{Archivo:} \texttt{test\_domain\_update\_delete\_performance.py}
  \item \textbf{Escenario:} Se crea un dominio semilla y se actualiza su campo
        \texttt{nombreDominio} mediante \texttt{Dominios.objects.filter(pk=...).update(...)}.
  \item \textbf{Verificación:} Se comprueba que se actualizó exactamente 1 registro, que el
        tiempo es inferior a 3.0~s y que el nuevo nombre existe en la base de datos.
  \item \textbf{Nota:} Marcada con TODO porque falta una vista/URL de actualización para validar
        el flujo HTTP completo.
\end{itemize}

\subsection{Aspectos positivos}

\begin{itemize}
  \item Los tiempos de ejecución son extremadamente bajos (del orden de milisegundos),
        lo que demuestra que las operaciones ORM son eficientes sobre SQLite.
  \item Las pruebas verifican tanto el resultado de la operación como el rendimiento temporal.
  \item El uso de \texttt{bulk\_create} para datos semilla permite un escenario realista
        con volúmenes significativos (200+ registros).
\end{itemize}

\subsection{Áreas de mejora}

\begin{itemize}
  \item Ambas pruebas están marcadas con \texttt{TODO}: se requiere implementar las
        vistas/URLs correspondientes para cubrir el flujo HTTP completo.
  \item Se recomienda agregar pruebas unitarias para validaciones de modelos, serialización
        de datos y lógica de negocio aislada de las vistas.
\end{itemize}

% ============================================================================
\newpage
\section{Pruebas de Integración}
% ============================================================================

Las pruebas de integración verifican la interacción correcta entre múltiples componentes
del sistema: vistas HTTP, modelos ORM, autenticación, y servicios externos simulados
mediante \texttt{monkeypatch}.

\subsection{Resultados}

\begin{center}
\rowcolors{2}{white}{secbg}
\begin{tabularx}{\textwidth}{>{\raggedright}p{3.5cm} >{\raggedright}p{3.2cm} r r >{\centering\arraybackslash}p{2.2cm}}
\toprule
\textbf{Prueba} & \textbf{Módulo} & \textbf{Tiempo (s)} & \textbf{Umbral (s)} & \textbf{Resultado} \\
\midrule
Disponibilidad de dominio
  & Dominios.verificar\_url
  & 0.014727
  & $< 1.0$
  & \aprobado \\
Creación de dominio y configuración
  & Dominios.agregar\_dominio
  & 0.048733
  & $< 3.0$
  & \aprobado \\
Listado de dominios
  & Clientes.mis\_hosts
  & 0.285564
  & $< 1.0$
  & \aprobado \\
Eliminación de dominio
  & Dominios.eliminar\_dominio
  & 0.013273
  & $< 3.0$
  & \aprobado \\
\bottomrule
\end{tabularx}
\end{center}

\subsection{Descripción de las pruebas}

\subsubsection{Consulta de disponibilidad de dominio}

\begin{itemize}
  \item \textbf{Archivo:} \texttt{test\_domain\_availability\_performance.py}
  \item \textbf{Escenario:} Se simula la verificación de disponibilidad del dominio
        \texttt{qa-disponible.test} mediante la vista \texttt{dominios:verificar\_url}.
        Se usa \texttt{monkeypatch} para simular fallos de red (\texttt{socket.gethostbyname}
        y \texttt{requests.get}).
  \item \textbf{Verificación:} Código HTTP $\neq$ 500, tiempo $< 1.0$~s, y presencia del
        nombre de dominio en la respuesta.
\end{itemize}

\subsubsection{Creación de dominio y configuración inicial}

\begin{itemize}
  \item \textbf{Archivo:} \texttt{test\_domain\_create\_configuration\_performance.py}
  \item \textbf{Escenario:} Se crea un dominio mediante POST a \texttt{dominios:agregar\_dominio}
        con un cliente autenticado. Se simulan los servicios de red y correo electrónico.
  \item \textbf{Verificación:} Redirección HTTP (302/303), tiempo $< 3.0$~s, y existencia
        del dominio en la base de datos.
\end{itemize}

\subsubsection{Listado de dominios del cliente}

\begin{itemize}
  \item \textbf{Archivo:} \texttt{test\_domain\_list\_performance.py}
  \item \textbf{Escenario:} Se crean 120 dominios semilla y se accede a la vista
        \texttt{clientes:mis\_hosts} con un cliente autenticado.
  \item \textbf{Verificación:} Código HTTP 200, tiempo $< 1.0$~s, y presencia del texto
        ``Mis Hosts'' en la respuesta.
  \item \textbf{Observación:} Es la prueba de integración más lenta (0.285~s), lo cual es
        esperado dado el volumen de datos renderizados.
\end{itemize}

\subsubsection{Eliminación de dominio}

\begin{itemize}
  \item \textbf{Archivo:} \texttt{test\_domain\_update\_delete\_performance.py}
  \item \textbf{Escenario:} Se crea un dominio semilla y se elimina mediante POST a
        \texttt{dominios:eliminar\_dominio}.
  \item \textbf{Verificación:} Redirección HTTP (302/303), tiempo $< 3.0$~s, y confirmación
        de que el registro ya no existe en la base de datos.
\end{itemize}

\subsection{Aspectos positivos}

\begin{itemize}
  \item Todas las pruebas de integración fueron aprobadas sin observaciones.
  \item El uso de \texttt{monkeypatch} para simular servicios externos (red, correo) permite
        ejecutar las pruebas de forma aislada sin dependencias externas.
  \item Se verifica tanto el código de respuesta HTTP como la persistencia en la base de datos,
        asegurando consistencia entre capas.
  \item El uso de fixtures (\texttt{cliente\_activo}, \texttt{cliente\_logueado}) promueve
        la reutilización y claridad del código de pruebas.
\end{itemize}

\subsection{Áreas de mejora}

\begin{itemize}
  \item Se recomienda agregar pruebas de integración para los módulos \texttt{Pagos},
        \texttt{Tickets}, \texttt{Administradores}, \texttt{Empleados} y \texttt{Distribuidor}.
  \item Incluir pruebas de autenticación y autorización (acceso denegado, roles incorrectos).
  \item Validar respuestas con contenido específico más allá de códigos HTTP.
\end{itemize}

% ============================================================================
\newpage
\section{Pruebas de Sistema}
% ============================================================================

Las pruebas de sistema verifican flujos completos de extremo a extremo, simulando
la experiencia real del usuario a través de múltiples módulos y operaciones secuenciales.

\subsection{Documentación de pruebas de sistema}

Las pruebas de sistema para el proyecto ChibchaWeb se encuentran documentadas en los
siguientes documentos adjuntos, los cuales definen el plan de calidad del software y los
resultados obtenidos tras su ejecución:

\begin{enumerate}
  \item \textbf{Plan de Calidad de Software}
        (\texttt{docs/plan\_calidad.tex}): Define los procesos, métricas, estándares,
        responsabilidades, herramientas y procedimientos de evaluación que aseguran la
        calidad del sistema ChibchaWeb durante todas sus fases. Incluye objetivos de
        calidad en seguridad, rendimiento, mantenibilidad, usabilidad y confiabilidad,
        así como las métricas cuantitativas y los roles del equipo.

  \item \textbf{Resultados del Plan de Calidad}
        (\texttt{docs/resultados\_plan\_calidad.tex}): Documenta los resultados iniciales
        de la medición de métricas de calidad tras el despliegue de la aplicación. Incluye
        los resultados de tiempo de respuesta del sistema (medido con Google PageSpeed),
        disponibilidad del sistema (medida con UptimeRobot), seguridad de los datos
        (auditoría con HostedScan) y satisfacción del usuario (encuestas con promedio
        de 4,76/5,0).
\end{enumerate}

Estos documentos cubren la planificación y verificación de la calidad del sistema a nivel
global, complementando las pruebas unitarias, de integración, de eficiencia y de robustez
documentadas en las secciones anteriores y posteriores de este informe.

% ============================================================================
\newpage
\section{Pruebas de Eficiencia}
% ============================================================================

Las pruebas de eficiencia verifican que el sistema cumple con los requisitos de rendimiento
en términos de tiempos de respuesta, throughput (transacciones por minuto), uso de recursos
computacionales y tamaño del proyecto.

\subsection{Resultados}

\begin{center}
\rowcolors{2}{white}{secbg}
\begin{tabularx}{\textwidth}{>{\raggedright}p{3cm} >{\raggedright}p{2.5cm} r r >{\centering\arraybackslash}p{2.8cm}}
\toprule
\textbf{Prueba} & \textbf{Tipo} & \textbf{Valor medido} & \textbf{Umbral} & \textbf{Resultado} \\
\midrule
Uso CPU/RAM bajo carga normal
  & Recursos
  & CPU: 0.00\% \newline RAM: 108.51 MB
  & CPU $\leq$ 70\% \newline RAM $\leq$ 1 GB
  & \aprobadojust \\
Throughput 200--500 TPM
  & Rendimiento
  & 36\,689.51 TPM
  & 200--500 TPM
  & \aprobadojust \\
Tamaño del proyecto
  & Almacenamiento
  & 0.71 MB
  & $< 250$ MB
  & \aprobado \\
\bottomrule
\end{tabularx}
\end{center}

\smallskip
{\footnotesize\textbf{*} \textit{APROBADO con justificación local}: el valor medido excede el rango
esperado en producción, pero se justifica por las condiciones del entorno de pruebas local
(SQLite en memoria, sin latencia de red, sin carga concurrente real).}

\subsection{Descripción de las pruebas}

\subsubsection{Uso de CPU y memoria bajo carga normal}

\begin{itemize}
  \item \textbf{Archivo:} \texttt{test\_resource\_usage\_performance.py}
  \item \textbf{Escenario:} Se crean 100 dominios semilla y se ejecutan 30 consultas
        \texttt{count()} consecutivas sobre el ORM. Se miden CPU y RAM mediante \texttt{psutil}.
  \item \textbf{Resultados:}
        \begin{itemize}
          \item CPU: 0.00\% (esperado en ejecución local breve)
          \item RAM: 108.51~MB (dentro del límite de 1~GB)
          \item Tiempo total: 0.063~s
        \end{itemize}
  \item \textbf{Justificación:} La RAM reportada (108.51~MB) es inferior al rango de referencia
        del acta (512~MB a 1~GB) porque la ejecución local con pytest no involucra un servidor
        WSGI permanente con múltiples workers.
\end{itemize}

\subsubsection{Throughput controlado (TPM)}

\begin{itemize}
  \item \textbf{Archivo:} \texttt{test\_throughput\_performance.py}
  \item \textbf{Escenario:} Se crean 80 dominios semilla y se ejecutan 20 transacciones de
        lectura secuenciales. Se calcula el TPM como
        $\text{TPM} = (\text{transacciones} / \text{duración\_s}) \times 60$.
  \item \textbf{Resultados:}
        \begin{itemize}
          \item TPM medido: 36\,689.51
          \item Duración total: 0.032~s
        \end{itemize}
  \item \textbf{Justificación:} El TPM supera ampliamente el rango esperado (200--500) porque
        la simulación local usa SQLite en memoria de proceso sin latencia de red ni usuarios
        reales concurrentes.
\end{itemize}

\subsubsection{Tamaño del proyecto comprimido}

\begin{itemize}
  \item \textbf{Archivo:} \texttt{test\_resource\_usage\_performance.py}
  \item \textbf{Escenario:} Se comprime todo el proyecto (excluyendo \texttt{.git},
        \texttt{venv}, \texttt{\_\_pycache\_\_} y cachés) en un archivo ZIP temporal.
  \item \textbf{Resultados:} 0.71~MB (muy por debajo del umbral de 250~MB).
  \item \textbf{Observación:} Demuestra que el proyecto mantiene un tamaño compacto y
        manejable sin dependencias pesadas incluidas en el repositorio.
\end{itemize}

\subsection{Resumen de tiempos de respuesta (todas las pruebas)}

La siguiente tabla consolida los tiempos de respuesta medidos en todas las pruebas de
la suite, ordenados de menor a mayor:

\begin{center}
\rowcolors{2}{white}{secbg}
\begin{tabularx}{\textwidth}{>{\raggedright}p{5.5cm} r r >{\centering\arraybackslash}X}
\toprule
\textbf{Escenario} & \textbf{Tiempo (s)} & \textbf{Umbral (s)} & \textbf{Margen} \\
\midrule
Actualización de dominio (ORM)           & 0.000862 & 3.0 & 99.97\% \\
Búsqueda/filtro de dominios (ORM)        & 0.002868 & 1.0 & 99.71\% \\
Eliminación de dominio (vista)           & 0.013273 & 3.0 & 99.56\% \\
Disponibilidad de dominio (vista)        & 0.014727 & 1.0 & 98.53\% \\
Throughput controlado (20 transacciones) & 0.032707 & 5.0 & 99.35\% \\
Alta carga concurrente (16 lecturas)     & 0.039506 & 5.0 & 99.21\% \\
Creación de dominio (vista)              & 0.048733 & 3.0 & 98.38\% \\
Uso CPU/RAM (30 consultas)               & 0.063045 & 5.0 & 98.74\% \\
Listado de dominios (vista, 120 registros) & 0.285564 & 1.0 & 71.44\% \\
Tamaño proyecto (compresión ZIP)         & 1.462848 & ---  & --- \\
\bottomrule
\end{tabularx}
\end{center}

\subsection{Aspectos positivos}

\begin{itemize}
  \item Todos los escenarios cumplen con los umbrales definidos con márgenes amplios
        (superiores al 70\% en todos los casos).
  \item La medición de recursos (CPU, RAM) mediante \texttt{psutil} proporciona evidencia
        cuantitativa objetiva.
  \item El cálculo de TPM permite estimar la capacidad de procesamiento del sistema.
  \item El control de tamaño del proyecto asegura que el repositorio no acumula artefactos
        innecesarios.
\end{itemize}

\subsection{Áreas de mejora}

\begin{itemize}
  \item Ejecutar pruebas de eficiencia con PostgreSQL para obtener métricas más representativas
        del entorno de producción.
  \item Implementar pruebas de carga con herramientas como \texttt{locust} o \texttt{k6} para
        simular múltiples usuarios concurrentes reales.
  \item Agregar pruebas de tiempo de respuesta para los módulos \texttt{Pagos}, \texttt{Tickets}
        y \texttt{Administradores}.
\end{itemize}

% ============================================================================
\newpage
\section{Pruebas de Robustez}
% ============================================================================

Las pruebas de robustez verifican la estabilidad del sistema bajo condiciones adversas,
incluyendo alta concurrencia, fallos de red simulados y cargas elevadas.

\subsection{Resultados}

\begin{center}
\rowcolors{2}{white}{secbg}
\begin{tabularx}{\textwidth}{>{\raggedright}p{3.8cm} >{\raggedright}p{3cm} r r >{\centering\arraybackslash}p{2.2cm}}
\toprule
\textbf{Prueba} & \textbf{Tipo} & \textbf{Tiempo (s)} & \textbf{Umbral (s)} & \textbf{Resultado} \\
\midrule
Alta carga/concurrencia de lectura
  & Concurrencia
  & 0.039506
  & $< 5.0$
  & \aprobado \\
Disponibilidad con fallo de red
  & Tolerancia a fallos
  & 0.014727
  & $< 1.0$
  & \aprobado \\
CPU/RAM bajo carga normal
  & Estabilidad de recursos
  & 0.063045
  & $< 5.0$
  & \aprobadojust \\
\bottomrule
\end{tabularx}
\end{center}

\subsection{Descripción de las pruebas}

\subsubsection{Alta carga y concurrencia de lectura}

\begin{itemize}
  \item \textbf{Archivo:} \texttt{test\_throughput\_performance.py}
  \item \textbf{Escenario:} Se crean 150 dominios semilla y se ejecutan 16 lecturas
        concurrentes mediante \texttt{ThreadPoolExecutor} con 4 workers.
  \item \textbf{Verificación:} Todas las lecturas retornan al menos 150 registros y el
        tiempo total es inferior a 5.0~s.
  \item \textbf{Análisis:} Demuestra que SQLite maneja correctamente lecturas concurrentes
        sin errores de bloqueo o corrupción de datos.
\end{itemize}

\subsubsection{Tolerancia a fallos de red}

\begin{itemize}
  \item \textbf{Archivo:} \texttt{test\_domain\_availability\_performance.py}
  \item \textbf{Escenario:} Se simula un fallo de red mediante \texttt{monkeypatch} sobre
        \texttt{socket.gethostbyname} y \texttt{requests.get}. La vista debe manejar
        graciosamente la excepción sin devolver HTTP~500.
  \item \textbf{Verificación:} La respuesta no es un error de servidor y contiene el nombre
        del dominio consultado, demostrando un manejo adecuado de excepciones.
\end{itemize}

\subsubsection{Estabilidad de recursos bajo carga}

\begin{itemize}
  \item \textbf{Archivo:} \texttt{test\_resource\_usage\_performance.py}
  \item \textbf{Escenario:} Se ejecutan 30 consultas consecutivas sobre 100 registros y se
        monitorean CPU y RAM para verificar que no hay fugas de memoria ni picos de CPU.
  \item \textbf{Verificación:} CPU $\leq$ 70\% y RAM $\leq$ 1~GB.
\end{itemize}

\subsection{Aspectos positivos}

\begin{itemize}
  \item La prueba de concurrencia demuestra que el sistema maneja múltiples lecturas
        simultáneas sin degradación ni errores.
  \item El manejo de fallos de red en la vista de disponibilidad muestra un diseño defensivo
        adecuado que no expone errores internos al usuario.
  \item El monitoreo de recursos proporciona evidencia de que el sistema no presenta fugas
        de memoria bajo carga repetida.
\end{itemize}

\subsection{Áreas de mejora}

\begin{itemize}
  \item Agregar pruebas de robustez con escrituras concurrentes (no solo lecturas).
  \item Implementar pruebas con datos corruptos o entradas malformadas para verificar
        la validación de datos.
  \item Agregar pruebas de recuperación (resiliencia): verificar que el sistema se recupera
        correctamente después de un fallo.
  \item Probar con un número mayor de workers concurrentes para identificar el punto de
        saturación.
\end{itemize}

% ============================================================================
\newpage
\section{Evidencia de ejecución}
% ============================================================================

\subsection{Salida completa de pytest}

\begin{lstlisting}[caption={Salida de ejecución de pytest},label={lst:pytest-output}]
============================= test session starts =============================
platform win32 -- Python 3.10.11, pytest-9.0.3, pluggy-1.6.0
django: version: 5.2.4, settings: ChibchaWeb.settings (from ini)
rootdir: C:\Users\Vortex\Documents\personal\ChibchaWeb
configfile: pytest.ini
plugins: anyio-4.12.0, django-4.12.0
collected 10 items

test/test_domain_availability_performance.py::
    test_domain_availability_view_under_one_second            PASSED [ 10%]
test/test_domain_create_configuration_performance.py::
    test_domain_create_and_initial_configuration_under_3s     PASSED [ 20%]
test/test_domain_list_performance.py::
    test_domain_list_view_under_one_second                    PASSED [ 30%]
test/test_domain_search_filter_performance.py::
    test_domain_search_filter_orm_under_one_second            PASSED [ 40%]
test/test_domain_update_delete_performance.py::
    test_domain_update_orm_under_three_seconds                PASSED [ 50%]
test/test_domain_update_delete_performance.py::
    test_domain_delete_view_under_three_seconds               PASSED [ 60%]
test/test_resource_usage_performance.py::
    test_cpu_and_memory_usage_under_normal_load               PASSED [ 70%]
test/test_throughput_performance.py::
    test_controlled_throughput_between_200_and_500_tpm        PASSED [ 80%]
test/test_throughput_performance.py::
    test_high_load_orm_scenario_under_five_seconds            PASSED [ 90%]
test/test_resource_usage_performance.py::
    test_compressed_project_size_under_250_mb                 PASSED [100%]

====================== 10 passed, 10 warnings in 38.33s =======================
\end{lstlisting}

\subsection{Comando de ejecución}

\begin{lstlisting}[language=bash,caption={Comando utilizado para ejecutar las pruebas}]
conda activate p
python -m pytest test/ -v --tb=long
\end{lstlisting}

\subsection{Configuración del entorno}

\begin{lstlisting}[caption={Archivo pytest.ini}]
[pytest]
DJANGO_SETTINGS_MODULE = ChibchaWeb.settings
pythonpath = .
\end{lstlisting}

% ============================================================================
\newpage
\section{Tabla consolidada de resultados}
% ============================================================================

\begin{center}
\rowcolors{2}{white}{secbg}
\small
\begin{longtable}{c >{\raggedright}p{3.2cm} >{\raggedright}p{2.2cm} r >{\raggedright}p{2cm} >{\centering\arraybackslash}p{2.2cm}}
\toprule
\textbf{No.} & \textbf{Escenario} & \textbf{Módulo} & \textbf{Tiempo (s)} & \textbf{Umbral} & \textbf{Resultado} \\
\midrule
\endhead
1 & Consulta de disponibilidad de dominio
  & Dominios.verificar\_url & 0.014727 & $< 1.0$ s & \aprobado \\
2 & Creación de dominio y configuración inicial
  & Dominios.agregar\_dominio & 0.048733 & $< 3.0$ s & \aprobado \\
3 & Listado de dominios del cliente
  & Clientes.mis\_hosts & 0.285564 & $< 1.0$ s & \aprobado \\
4 & Búsqueda/filtro de dominios
  & Dominios ORM & 0.002868 & $< 1.0$ s & \aprobadotodo \\
5 & Actualización de dominio
  & Dominios ORM & 0.000862 & $< 3.0$ s & \aprobadotodo \\
6 & Eliminación de dominio
  & Dominios.eliminar\_dominio & 0.013273 & $< 3.0$ s & \aprobado \\
7 & Uso CPU/RAM bajo carga normal
  & Proceso pytest & 0.063045 & CPU $\leq$ 70\% & \aprobadojust \\
8 & Throughput 200--500 TPM
  & Dominios ORM & 0.032707 & 200--500 TPM & \aprobadojust \\
9 & Alta carga/concurrencia
  & Dominios ORM & 0.039506 & $< 5.0$ s & \aprobado \\
10 & Tamaño proyecto comprimido
  & Repositorio & 1.462848 & $< 250$ MB & \aprobado \\
\bottomrule
\end{longtable}
\end{center}

% ============================================================================
\section{Conclusiones}
% ============================================================================

\begin{enumerate}
  \item \textbf{Resultado general:} Las 10 pruebas de la suite fueron aprobadas exitosamente.
        No se registraron fallos ni errores de ejecución.

  \item \textbf{Rendimiento:} Todos los escenarios cumplen holgadamente con los umbrales de
        tiempo definidos. El margen mínimo observado es del 71.44\% (listado de 120 dominios),
        lo cual sigue siendo ampliamente satisfactorio.

  \item \textbf{Cobertura por categoría:}
        \begin{itemize}
          \item \textit{Unitarias:} 2 pruebas — cobertura básica sobre capa ORM.
          \item \textit{Integración:} 4 pruebas — buena cobertura de vistas HTTP principales.
          \item \textit{Sistema:} documentadas en \texttt{plan\_calidad.tex} y \texttt{resultados\_plan\_calidad.tex}.
          \item \textit{Eficiencia:} 3 pruebas — cobertura de CPU, RAM, TPM y tamaño.
          \item \textit{Robustez:} 3 pruebas — cobertura de concurrencia y tolerancia a fallos.
        \end{itemize}

  \item \textbf{Observaciones:}
        \begin{itemize}
          \item Dos pruebas están marcadas con \texttt{TODO} pendiente de implementación de
                vistas/URLs para completar el flujo HTTP.
          \item Dos pruebas fueron aprobadas con justificación local debido a que las métricas
                de TPM y RAM difieren del rango esperado en producción.
          \item Las 10 advertencias registradas corresponden a la marca \texttt{@pytest.mark.performance}
                no registrada formalmente en la configuración de pytest.
        \end{itemize}

  \item \textbf{Recomendaciones:}
        \begin{itemize}
          \item Continuar ampliando las pruebas de sistema (E2E) para los flujos principales.
          \item Ejecutar la suite con PostgreSQL para obtener métricas representativas de producción.
          \item Ampliar la cobertura a los módulos \texttt{Pagos}, \texttt{Tickets},
                \texttt{Administradores} y \texttt{Distribuidor}.
          \item Registrar la marca \texttt{performance} en \texttt{pytest.ini} para eliminar
                las advertencias.
        \end{itemize}
\end{enumerate}

\end{document}
